\begin{abstract}
배틀로얄 게임의 공개 텔레메트리는 줄어드는 활동 영역에서 다수의 팀이 경쟁하는 과정을 대규모 시공간 데이터로 기록한다.
본 연구는 단계 경계에서 생존한 팀이 다음 단계 안에 탈락할 조건부 확률을 서식지 축소 관점의 단계 조건부 이산시간 hazard로 정식화하고, 이 위험의 예측을 무엇이 좌우하는지 분석한다.
PUBG 13,948경기의 텔레메트리 약 3억 2천만 건에서 예측 시점 이후의 정보가 섞이지 않는 팀--단계 표본 817,766개를 구성하고, 12개 코호트에서 활동 점수 기반 설계와 상대위험 기반 설계의 적용판을 같은 조건으로 비교한 뒤 설계 요소를 하나씩 바꾸어 차이를 분해하였다.
주변 경쟁 밀도와 이동 부담은 위험을 높이고 남은 팀원 수는 낮추었으며, 구역 안쪽 깊이의 보호 효과는 구역이 작아질수록 커졌고, 해석 가능한 네 변수가 고차원 모델 판별력의 약 79\%를 설명하였다.
본 연구의 설계는 모든 코호트에서 두 적용판보다 높은 AUC(0.778 대 0.686--0.763)를 보였고, 상대위험 기반 적용판이 거의 평평하게 예측한 단계별 위험 수준을 정확히 나타냈다(단계별 보정 오차 0.013 대 0.064).
활동 점수 기반 설계와의 차이는 주로 서식지 맥락이, 상대위험 기반 설계와의 차이는 주로 단계별 기저위험과 학습 목적이 설명하였으며, 모델 구조를 바꾼 이득은 AUC 0.005 미만이었다.
\end{abstract}

% 주제어는 EasyChair 제출 양식에 입력(본문에서는 분량 때문에 생략):
% 시공간 텔레메트리, 이산시간 hazard, 데이터 누수, 생존분석, e스포츠 분석
